% Shared setup of the regression tests: `\input stridi-test.tex\relax` at the top of each test.
%
% The package is loaded here, before \START, so that its version line never reaches the part of the log compared with the .tlg files: `l3build ctan` runs the tests after `l3build tag` has stamped the version.
% expl3's debug checks are on, so the tests also catch misused variables and deprecated functions.
\input regression-test.tex\relax
\documentclass{article}
\RequirePackage[enable-debug]{expl3}
\ExplSyntaxOn
\debug_on:n { check-declarations , check-expressions , deprecation }
\ExplSyntaxOff
\RequirePackage{string-diagrams}

% Use Computer Modern on all engines: XeTeX and LuaTeX default to Latin Modern, whose slightly different em would change every size based on it.
\renewcommand\encodingdefault{OT1}
\renewcommand\rmdefault{cmr}

% Load the math fonts before \START, so the font messages of the first box label don't reach the logs.
\AtBeginDocument{\sbox0{$x$}}

% NOTE: TikZ remembers node names across pictures, so a test about a missing node must use a name that no other test defines.

% \PICTURE[<options>]{<code>}: typeset a tikzpicture in a box that is never shipped out.
\NewDocumentCommand\PICTURE{O{}+m}{%
  \setbox0\hbox{\begin{tikzpicture}[#1]#2\end{tikzpicture}}%
}

\makeatletter

% \LOG{<code>}: in a picture, log the PDF operators that <code> draws with, leaving out node labels.
% NOTE: each driver places labels with its own operators, so they're dropped to keep the logs the same on all engines.
% NOTE: colors are set through the color stack rather than PDF operators, so they don't show.
\protected\long\def\LOG#1{%
  \begingroup
    \def\pgfsys@invoke##1{\TYPE{##1}}%
    \def\pgfsys@hbox##1{}%
    \def\pgfsys@hboxsynced##1{}%
    #1%
  \endgroup
}

% \SHOWPOINT{<label>}{<pgf point>}: log the coordinates of a point.
\protected\long\def\SHOWPOINT#1#2{%
  #2%
  \TYPE{#1 = (\the\pgf@x, \the\pgf@y)}%
}

% \SHOWANCHOR{<node>}{<anchor>}: log the coordinates of an anchor.
\protected\def\SHOWANCHOR#1#2{%
  \SHOWPOINT{#1.#2}{\pgfpointanchor{#1}{#2}}%
}

% \SHOWCOORD{<coordinate>}: log the point that TikZ makes of a coordinate.
\protected\def\SHOWCOORD#1{%
  \SHOWPOINT{(#1)}{\tikz@scan@one@point\pgfutil@firstofone(#1)}%
}

\makeatother

% A box size that gives round coordinates.
\tikzset{
  square/.style={minimum size=120pt, inner sep=0pt, outer sep=0pt},
}
